\begin{figure}[htbp]
\centering
\begin{tikzpicture}[x=.95cm,y=.95cm,>=Latex,font=\small,
  primal/.style={draw=blue!65!black,very thick},
  dual/.style={draw=orange!80!black,thick,densely dashed}]
  \path[draw=black!40,fill=blue!6]
    (0,-.69282)--(0,.69282)--(-1.2,1.38564)--(-2.4,.69282)
    --(-2.4,-.69282)--(-1.2,-1.38564)--cycle;
  \path[draw=black!40,fill=orange!7]
    (0,-.69282)--(1.2,-1.38564)--(2.4,-.69282)--(2.4,.69282)
    --(1.2,1.38564)--(0,.69282)--cycle;
  \draw[primal,->] (0,-.69282)--(0,.69282);
  \node[anchor=south] at (0,.90) {$e$};
  \draw[dual,->] (1.2,0)--(-1.2,0);
  \fill (-1.2,0) circle (1.8pt) node[below=4pt] {$y$};
  \fill (1.2,0) circle (1.8pt) node[below=4pt] {$z$};
  \node at (-1.2,.68) {left face};
  \node at (1.2,.68) {right face};
  \node at (0,-1.67) {crossing sign $+1$};
  \node at (0,-2.3) {(a) Dual edge $[y,z]$, traversed from $z$ to $y$};

  \begin{scope}[xshift=6.65cm]
    \path[fill=blue!6]
      (1.5,-.7)--(1.5,.7)--(0,1.55)--(-1.5,.7)
      --(-1.5,-.7)--(0,-1.55)--cycle;
    \draw[primal]
      (1.5,-.7)--(1.5,.7)--(0,1.55)--(-1.5,.7)
      --(-1.5,-.7)--(0,-1.55)--cycle;
    \draw[primal,->] (1.5,-.45)--(1.5,.5);
    \draw[primal,->] (-.33,1.363)--(-.95,1.012);
    \node[anchor=south] at (0,1.82) {$L$ (counterclockwise)};
    \draw[black!45] (2.7,-1.65)--(2.7,1.65)
      node[above] {$\partial D$};
    \draw[dual,->] (2.7,0)--(0,0);
    \fill (2.7,0) circle (1.8pt) node[right=3pt] {$y_0$};
    \fill (0,0) circle (1.8pt) node[below=4pt] {$y$};
    \node at (0,-.78) {$\operatorname{cr}_\Gamma(L)=1$};
    \node at (.4,-2.3) {(b) Net crossing equals enclosure};
  \end{scope}
\end{tikzpicture}
\caption{The crossing convention for the current primitive.
(a) The oriented honeycomb link $e$ crosses the tiling edge $[y,z]$;
the dual traversal from its right face to its left face has sign $+1$.
Reversing either orientation reverses this sign.
(b) For a counterclockwise contractible polygon, the interior lies to
the left. Hence a dual chain from the exterior boundary gap $y_0$ to
$y$ has net crossing $1$ if the polygon encloses $y$, and $0$ otherwise.
The global polygon and chain are schematic. These are topological
crossing signs; the imaginary prefix-current coefficient is computed
separately in the proof.}
\label{fig:geometry-current}
\end{figure}
